% GENERATED FILE. Edit canonical lessons, then run npm run generate:books.
% canonical-source-hash: fnv1a64:541da917af2abb2f
% canonical-lessons: PT-C459-rochedo, PT-C459-mare, PT-C459-concha, PT-C459-bronzear, PT-C459-mergulhar

\chapter{Nas rochas}
\label{ch:pt-nas-rochas}
\hlchaptermodality{459}

\begin{chapteropening}
\textbf{By the end of this chapter:} \emph{I can describe a day on a rocky shore: the rocks, the tide, seashells, getting a tan and diving in.}

Complete the last lesson of chapter 459: I can describe a day on a rocky shore: the rocks, the tide, seashells, getting a tan and diving in.
\end{chapteropening}

\section[o rochedo]{o rochedo — the big rock, the crag}

\label{lesson:PT-C459-rochedo}



\begin{quote}
Before the new one: say the Portuguese for the gorge, then the Portuguese for to climb.
\end{quote}

\subsection*{You'll want to know: o rochedo}
\textbf{o rochedo} — \textquotedblleft{}the big rock, the crag\textquotedblright{}.

\begin{quote}
Sentámo-nos num rochedo a ver o mar.
\end{quote}

\begin{quote}
\emph{We sat on a big rock watching the sea.}
\end{quote}

Built on \emph{a rocha}.

\begin{cousinweb}
\emph{rochedo} is built on \emph{rocha} 'rock' with \emph{-edo}, a group of things. \emph{Rocha} came from French \emph{roche} (EtymDB), whose earlier origin is not settled.
\end{cousinweb}

\begin{grammarlens}[title={What you've built}]
One more word for telling, explaining and describing at B1.
\end{grammarlens}

\subsection*{Your turn}
\begin{itemize}
\raggedright
  \item \emph{Say it:} \emph{o rochedo}
  \item \emph{Say it:} \emph{o rochedo}, once more
  \item \emph{Say it:} a sentence of your own with \emph{o rochedo}
  \item \emph{Recall:} say the Portuguese for the gorge, then the Portuguese for to climb, then say \emph{o rochedo} again
\end{itemize}

\begin{tcolorbox}[breakable,colback=teal!4,colframe=teal!35!black,title={Before you move on}]
What does o rochedo mean? (\textquotedblleft{}The big rock, the crag\textquotedblright{}.) Say it once more.
\end{tcolorbox}
\section[a maré]{a maré — the tide}

\label{lesson:PT-C459-mare}



\begin{quote}
Before the new one: say the Portuguese for to climb, then the Portuguese for the big rock.
\end{quote}

\subsection*{You'll want to know: a maré}
\textbf{a maré} — \textquotedblleft{}the tide\textquotedblright{}.

\begin{quote}
Com a maré baixa, podemos andar até à ilha.
\end{quote}

\begin{quote}
\emph{At low tide, we can walk to the island.}
\end{quote}

\emph{Maré alta} — high tide; \emph{maré baixa} — low tide.

\begin{cousinweb}
\emph{maré} came from French \emph{marée} 'tide', built on \emph{mer} 'sea', from Latin \emph{mare} (EtymDB).
\end{cousinweb}

\begin{grammarlens}[title={What you've built}]
One more word for telling, explaining and describing at B1.
\end{grammarlens}

\subsection*{Your turn}
\begin{itemize}
\raggedright
  \item \emph{Say it:} \emph{a maré}
  \item \emph{Say it:} \emph{a maré}, once more
  \item \emph{Say it:} a sentence of your own with \emph{a maré}
  \item \emph{Recall:} say the Portuguese for to climb, then the Portuguese for the big rock, then say \emph{a maré} again
\end{itemize}

\begin{tcolorbox}[breakable,colback=teal!4,colframe=teal!35!black,title={Before you move on}]
What does a maré mean? (\textquotedblleft{}The tide\textquotedblright{}.) Say it once more.
\end{tcolorbox}
\section[a concha]{a concha — the seashell; the ladle}

\label{lesson:PT-C459-concha}



\begin{quote}
Before the new one: say the Portuguese for the big rock, then the Portuguese for the tide.
\end{quote}

\subsection*{You'll want to know: a concha}
\textbf{a concha} — \textquotedblleft{}the seashell; the ladle\textquotedblright{}.

\begin{quote}
A minha filha apanhou conchas na praia.
\end{quote}

\begin{quote}
\emph{My daughter picked up shells on the beach.}
\end{quote}

In the kitchen, \emph{a concha} is the soup ladle.

\begin{cousinweb}
\emph{concha} is inherited from Vulgar Latin \emph{concla}, from Latin \emph{conchula}, a small form of \emph{concha} 'shell' (EtymDB).
\end{cousinweb}

\begin{grammarlens}[title={What you've built}]
One more word for telling, explaining and describing at B1.
\end{grammarlens}

\subsection*{Your turn}
\begin{itemize}
\raggedright
  \item \emph{Say it:} \emph{a concha}
  \item \emph{Say it:} \emph{a concha}, once more
  \item \emph{Say it:} a sentence of your own with \emph{a concha}
  \item \emph{Recall:} say the Portuguese for the big rock, then the Portuguese for the tide, then say \emph{a concha} again
\end{itemize}

\begin{tcolorbox}[breakable,colback=teal!4,colframe=teal!35!black,title={Before you move on}]
What does a concha mean? (\textquotedblleft{}The seashell; the ladle\textquotedblright{}.) Say it once more.
\end{tcolorbox}
\section[bronzear]{bronzear — to tan, to get a tan}

\label{lesson:PT-C459-bronzear}



\begin{quote}
Before the new one: say the Portuguese for the tide, then the Portuguese for the seashell.
\end{quote}

\subsection*{You'll want to know: bronzear}
\textbf{bronzear} — \textquotedblleft{}to tan, to get a tan\textquotedblright{}.

\begin{quote}
Passámos a tarde a bronzear-nos na areia.
\end{quote}

\begin{quote}
\emph{We spent the afternoon tanning on the sand.}
\end{quote}

\emph{Bronzear-se}; \emph{estar bronzeado} — to be tanned.

\begin{cousinweb}
\emph{bronzear} is built on \emph{bronze}, which was borrowed from French \emph{bronze}, from Italian \emph{bronzo} (EtymDB): skin the colour of bronze.
\end{cousinweb}

\begin{grammarlens}[title={What you've built}]
One more word for telling, explaining and describing at B1.
\end{grammarlens}

\subsection*{Your turn}
\begin{itemize}
\raggedright
  \item \emph{Say it:} \emph{bronzear}
  \item \emph{Say it:} \emph{bronzear}, once more
  \item \emph{Say it:} a sentence of your own with \emph{bronzear}
  \item \emph{Recall:} say the Portuguese for the tide, then the Portuguese for the seashell, then say \emph{bronzear} again
\end{itemize}

\begin{tcolorbox}[breakable,colback=teal!4,colframe=teal!35!black,title={Before you move on}]
What does bronzear mean? (\textquotedblleft{}To tan, to get a tan\textquotedblright{}.) Say it once more.
\end{tcolorbox}
\section[mergulhar]{mergulhar — to dive (in), to plunge}

\label{lesson:PT-C459-mergulhar}



\begin{quote}
Before the new one: say the Portuguese for the seashell, then the Portuguese for to tan.
\end{quote}

\subsection*{You'll want to know: mergulhar}
\textbf{mergulhar} — \textquotedblleft{}to dive (in), to plunge\textquotedblright{}.

\begin{quote}
Ele mergulhou do rochedo para a água.
\end{quote}

\begin{quote}
\emph{He dived off the rock into the water.}
\end{quote}

Also to dip something: \emph{mergulhar o pão no café}.

\begin{cousinweb}
\emph{mergulhar} is inherited from Vulgar Latin \emph{merguliāre} (EtymDB), from \emph{mergulus} 'little diver (bird)', on \emph{mergere} 'to dip'.
\end{cousinweb}

\begin{grammarlens}[title={What you've built}]
One more word for telling, explaining and describing at B1.
\end{grammarlens}

\subsection*{Your turn}
\begin{itemize}
\raggedright
  \item \emph{Say it:} \emph{mergulhar}
  \item \emph{Say it:} \emph{mergulhar}, once more
  \item \emph{Say it:} a sentence of your own with \emph{mergulhar}
  \item \emph{Recall:} say the Portuguese for the seashell, then the Portuguese for to tan, then say \emph{mergulhar} again
\end{itemize}

\begin{tcolorbox}[breakable,colback=teal!4,colframe=teal!35!black,title={Before you move on}]
What does mergulhar mean? (\textquotedblleft{}To dive (in), to plunge\textquotedblright{}.) Say it once more.
\end{tcolorbox}
